\section{Table examples}

Use a compact table to report exact values that readers may want to compare.
Give units in the heading, align comparable numbers, and put the caption
above the table. Table~\ref{tab:example-values} shows the first and final
values from the example plot.

\begin{table}[htbp]
  \centering
  \caption{Illustrative response times at the first and final trials.}
  \label{tab:example-values}
  \begin{tabular}{lr}
    \toprule
    Trial & Response time (ms) \\
    \midrule
    First & 340 \\
    Final & 300 \\
    \bottomrule
  \end{tabular}
\end{table}

When a column needs longer descriptions, use a wrapping column.
Table~\ref{tab:example-measures} shows how to define measures and give
enough context for a reader to interpret them.

\begin{table}[htbp]
  \centering
  \caption{Example definitions for measures reported in a user study.}
  \label{tab:example-measures}
  \begin{tabularx}{\linewidth}{lX}
    \toprule
    Measure & Definition \\
    \midrule
    Response time & Elapsed time from the appearance of the prompt to the
      participant's first recorded response, in milliseconds. \\
    Completion rate & Percentage of trials in which the participant completed
      the task within the time limit. \\
    \bottomrule
  \end{tabularx}
\end{table}
